\honest{What Evidence Filtered Evidence?}{Eliezer Yudkowsky}{2007}

Back to the clever arguer from ``The Bottom Line.'' Suppose the arguer may say only true things, and points out that box B has a blue stamp, which diamond boxes are more likely to have. ``Each statement that the clever arguer makes is valid evidence,'' and Jaynes says a Bayesian must condition on all known evidence, so it seems you must update. But then an arguer with enough signs to choose from could make you believe anything. ``That doesn't sound right.'' \nb{This is the question ``The Bottom Line'' left open, and this post answers it.}

Take a simpler case. A coin is biased either two-thirds heads or two-thirds tails, and each head is ``1 bit of evidence'' for the first. \nb{The footnote defines a bit as the logarithm of a probability. The ``1 bit'' here is the logarithm of a likelihood ratio, 2 to 1, a different quantity. ``How Much Evidence Does It Take?'' makes the same switch.} I flip it ten times and tell you that flips 4, 6 and 9 came up heads. What should you believe? It depends on my rule for choosing what to tell you. If I always report those three flips, the odds are 8 to 1 for heads. If I report every head and only heads, the other seven were tails, and the odds are 1 to 16. I might also have decided to report those three flips only if the coin was very likely heads-biased; and so on.

The Monty Hall problem is the same. You pick door 1 of three, one of which hides \$100,000, and the host opens door 2, which is empty. If the host always opens an empty door, you should switch to door 3. If the host always opens door 2 whatever is behind it, doors 1 and 3 are equally likely. If the host opens a door only when you picked the money, you should stay. Many people get the standard problem wrong because they update only on door 2 being empty. You must condition not only on door 2 being empty but on the host's choosing to open it. In general, you condition on a speaker with some rule having said the words, not only on what the words say. \nb{In the coin and door cases the listener knows the rule and the number of flips or doors. The post does not work through its opening case, in which the listener does not know how many signs the arguer left out. Yudkowsky worked it out in a comment the next day.}

``Most legal processes'' rely on two opposed advocates, since ``it is easier to find two biased humans than one unbiased one.'' With two boxes that is almost as good as one curious inquirer, but real questions often have many sides. \nb{Two opposed advocates are the theory of common-law courts. In the inquisitorial systems of continental Europe, Latin America and East Asia, the court investigates the facts itself.}

Do not use filtering as an excuse to dismiss evidence you dislike. ``If you're ticked off by a contrary argument, then you are familiar with the case,'' and you probably know your own side's best arguments already. So ``a blue stamp on box B is still evidence.'' \nb{Being annoyed shows that you care about a question, not that you know it.} Only if you hear one side for the first time should you beware. In a way, no one can really trust natural selection until they have listened to creationists for five minutes.
